%% kapa: exam=1 round=33 year=2022 period=2 form=A subject=accounting subject_ko=회계학 seq=20 number=60
%% source: Q-Net 감정평가사 기출문제 (cst00309) artlSeq=5209997 file=raw/1008/5209997/2022년도 제33회 감정평가사 1차시험 2교시(감정평가 관계법규, 회계학).pdf
%% conversion: kordoc 4.24.51 -> tools/gen1.py (auto); review=proofread-sonnet; flags=none
\begin{question}{60}
(주)감평은 20×1년 1월 1일에 액면금액 ￦500,000(표시이자율 연 10\%, 만기 3년, 매년 말 이자 지급)의 사채를 ￦475,982에 취득하고, 당기손익-공정가치 측정 금융자산으로 분류하였다. 동 사채의 취득 당시 유효이자율은 연 12\%이며, 20×1년 말 공정가치는 ￦510,000이다. 상기 금융자산(사채) 관련 회계처리가 (주)감평의 20×1년도 당기순이익에 미치는 영향은? (단, 단수차이로 인한 오차가 있다면 가장 근사치를 선택한다.)
\begin{choices}
  \item ￦84,018 증가
  \item ￦70,000 증가
  \item ￦60,000 증가
  \item ￦34,018 증가
  \item ￦10,000 증가
\end{choices}
\end{question}
\answer{1}
